\section{Further Robustness}
\label{app:online}

This appendix states sufficient conditions for the two-country sign result of \cref{sec:bilateral}. Throughout, $e_{Aj}$ is units of $A$'s currency per unit of $j$'s currency, so a rise is a depreciation of $A$ against $j$.

\subsection{Two-country sufficient conditions}
\label{app:online-twocountry}

Let $e\equiv e_{AB}$ and let $\tau$ be $A$'s tariff on imports from $B$. Border prices are one in the producer's currency, so with import and export quantities $Q_M$ and $Q_X$, $A$'s import payments are $M_A=eQ_M$, its export receipts are $X_A=Q_X$, and $\TB_A(e,\tau)\equiv X_A(e,\tau)-M_A(e,\tau)$, where $X_A$ and $M_A$ include all induced effects, including income effects of tariff revenue. If $\TB_A$ is continuously differentiable with $\TB_{A,e}\neq0$ at an equilibrium, the implicit function theorem gives a locally unique $e(\tau)$ with $de/d\tau=-\TB_{A,\tau}/\TB_{A,e}$.

The first condition is $\TB_{A,\tau}>0$. Because $M_A=eQ_M$ is measured at border prices, import payments fall at fixed $e$ if and only if the import quantity falls. Under a lump-sum rebate of revenue, with $\varepsilon^H$ the compensated own-price elasticity of import demand, $\eta_I$ its income elasticity and $\theta$ the revenue share of income,
\begin{equation*}
\frac{d\log Q_M}{d\log(1+\tau)}\bigg|_e=-\frac{\varepsilon^H}{1-\eta_I\theta},
\end{equation*}
so at free trade ($\theta=0$) the import quantity falls if and only if $\varepsilon^H>0$. In the nested-CES model at free trade, $\varepsilon^H=\rho_D>0$, the domestic diversion margin of \cref{eq:twocountry}, and $X_{A,\tau}=0$ because $B$'s income and demand are unchanged at a fixed exchange rate. The second condition is $\TB_{A,e}>0$, which at a balanced baseline is equivalent to the Marshall--Lerner condition $\varepsilon_X+\varepsilon_M>1$ (\cref{sec:bilateral}).

When both conditions hold, $de/d\tau<0$. These are sufficient conditions, not properties of every model: the fiscal closure (revenue wasted, spent by the government or transferred abroad) and the pricing convention (local- or dominant-currency pricing, incomplete pass-through, markups) change the quantity responses and can reverse either sign.

The online supplement also considers sector-specific factors and permanent trade imbalances. These extensions preserve the qualitative mechanism while shifting the quantitative bystander threshold.
